\documentclass[a4paper,12pt]{article}
\usepackage{color}
\usepackage{graphicx, gensymb}
\usepackage{amsfonts, cancel, amssymb}
\usepackage{amsmath, amsthm, array, esvect, esint}
\usepackage{typearea, tikz, pgffor, verbatim}
\usetikzlibrary{calc, quotes}
\usepackage[a4paper,left=2cm,right=2cm,top=2cm,bottom=3cm,bindingoffset=5mm]{geometry}
\newcommand\numberthis{\addtocounter{equation}{1}\tag{\theequation}}
\newcommand{\bra}{}
\newcommand{\kom}{}
\newcommand{\brak}{}
\newcommand{\braket}{}
\newcommand{\intx}{}
\newcommand{\intr}{}
\newcommand{\kla}{}
\newcommand{\klam}{}
\newcommand{\klammer}{}
\newcommand{\oper}{}
\newcommand{\tecxt}{}
\newcommand{\Bra}{}
\newcommand{\Ket}{}

\begin{document}

\title{02 Koordinatensysteme}
\author{Joachim Schwardt}
\date{\today}
\maketitle

\section*{Aufgabe 4: Nicht-orthogonale Koordinatensysteme}
Gegeben ist $\vec{r}=\binom{x^1+x^2}{x^2}$.
\subsection*{a)}Berechnung von $\vec{g}_a$ und $\vec{g}^a$:
\begin{align*}
\vec{g}_1&=\partial_1\vec{r}=\binom{1}{0}&\vec{g}_2&=\partial_2\vec{r}=\binom{1}{1}\numberthis\label{eqn:g unten a}\\
\vec{g}^1&=g^{1b}\vec{g}_b=\binom{1}{-1}&\vec{g}^2&=g^{2b}\vec{g}_b=\binom{0}{1}\numberthis\label{eqn:g oben a}
\end{align*}
Alternativ kann auch die Orthogonalitätsbedingung $\vec{g}_a\cdot\vec{g}^b=\delta_a^b$ verwendet werden. Es sei $\vec{g}^{1}=\binom{\alpha_{1}}{\beta_{1}}$ und $\vec{g}^{2}=\binom{\alpha_{2}}{\beta_{2}}$:
\begin{align*}
1&\overset{!}{=}\vec{g}_1\cdot\vec{g}^1=\alpha_1&0&\overset{!}{=}\vec{g}_2\cdot\vec{g}^1=\alpha_1+\beta_1\ \ \Rightarrow\ \ \beta_1=-1\\
0&\overset{!}{=}\vec{g}_1\cdot\vec{g}^2=\alpha_2&1&\overset{!}{=}\vec{g}_2\cdot\vec{g}^2=\alpha_2+\beta_2\ \ \Rightarrow\ \ \beta_2=1
\end{align*}
Graphische Darstellung der Vektoren in (\ref{figure:g oben unten a}):
\begin{figure}[h!]
\centering
\begin{tikzpicture}[scale=2]
\draw[->, >=latex] (0,-1) -- (0,2) node[right]{$x^2$};
\draw[->, >=latex] (-2,0) -- (2,0) node[below]{$x^1$};
\draw[->, >=latex, blue, line width=2] (0,0) -- (1,0) node[above]{$\vec{g}_1$};
\draw[->, >=latex, blue, line width=2] (0,0) -- (1,1) node[above]{$\vec{g}_2$};
\draw[->, >=latex, red, line width=2] (0,0) -- (1,-1) node[right]{$\vec{g}^1$};
\draw[->, >=latex, red, line width=2] (0,0) -- (0,1) node[left]{$\vec{g}^2$};
\end{tikzpicture}
\caption{Darstellung der Vektoren $\vec{g}_a$ und $\vec{g}^a$}
\label{figure:g oben unten a}
\end{figure}
\subsection*{b)}Die Metrik und ihre Inverse sind:
\begin{align*}
g_{ab}&=\begin{pmatrix}
1&1 \\ 1&2
\end{pmatrix}&\Rightarrow&&g^{ab}&=\begin{pmatrix}
2&-1 \\ -1&1
\end{pmatrix}\numberthis\label{eqn:Metrik}
\end{align*}
\subsection*{c)}Berechung der kinetischen Energie in Abhängigkeit von $\dot{x}^a$:
\begin{align*}
T&=\frac{m}{2}g_{ab}\dot{x}^a\dot{x}^b=\frac{m}{2}\kla\left[ {(\dot{x}^1)^2+2\dot{x}^1\dot{x}^2+2(\dot{x}^2)^2} \right]\numberthis\label{eqn:Kinetische Energie von xdot oben a}
\end{align*}
\subsection*{d)}Bestimmung des Gradienten einer Funktion $f(x^1,x^2)$:
\begin{align*}
\nabla f&=\vec{g}^a\partial_a f=\binom{1}{-1}\partial_1f+\binom{0}{1}\partial_2f=\binom{\partial_1f}{-\partial_1f+\partial_2f}\numberthis\label{eqn:Gradient}
\end{align*}
\subsection*{e)}Berechnung der Koeffizienten für $\vec{A}=\binom{1}{0}$ und $\vec{B}=\binom{0}{1}$:
\begin{align*}
A_1&=\vec{A}\cdot\vec{g}_1=1&A_2&=\vec{A}\cdot\vec{g}_2=1\\
A^1&=\vec{A}\cdot\vec{g}^1=1&A^2&=\vec{A}\cdot\vec{g}^2=0\\
B_1&=\vec{B}\cdot\vec{g}_1=0&B_2&=\vec{B}\cdot\vec{g}_2=1\\
B^1&=\vec{B}\cdot\vec{g}^1=-1&B^2&=\vec{B}\cdot\vec{g}^2=1
\end{align*}
\section*{Aufgabe 5: Zylinderkoordinaten}
Gegeben ist $\vec{r}=\begin{pmatrix}
x^1\cos x^2 \\ x^1\sin x^2 \\ x^3
\end{pmatrix}$ ($x^1=\rho$, $x^2=\phi$ und $x^3=z$).
\subsection*{a)}Berechnung von $\vec{g}_a$ und $\vec{g}^a$:
\begin{align*}
\vec{g}_1&=\partial_1\vec{r}=\begin{pmatrix}
\cos x^2 \\ \sin x^2 \\ 0
\end{pmatrix} &\vec{g}_2&=\partial_2\vec{r}=\begin{pmatrix}
-x^1\sin x^2 \\ x^1\cos x^2 \\ 0
\end{pmatrix}&\vec{g}_3&=\partial_3\vec{r}=\begin{pmatrix}
0 \\ 0 \\ 1
\end{pmatrix} \numberthis\label{eqn:zylinder:g unten a}\\
\vec{g}^1&=g^{1b}\vec{g}_b=\begin{pmatrix}
\cos x^2 \\ \sin x^2 \\ 0
\end{pmatrix}&\vec{g}^2&=g^{2b}\vec{g}_b=\begin{pmatrix}
-\frac{1}{x^1}\sin x^2 \\ \frac{1}{x^1}\cos x^2 \\ 0
\end{pmatrix}&\vec{g}^3&=g^{3b}\vec{g}_b=\begin{pmatrix}
0 \\ 0 \\ 1
\end{pmatrix}\numberthis\label{eqn:zylinder:g oben a}
\end{align*}
\subsection*{b)}Die Metrik und ihre Inverse sind:
\begin{align*}
g_{ab}&=\begin{pmatrix}
1&0&0 \\ 0&(x^1)^2&0 \\ 0&0&1
\end{pmatrix}&\Rightarrow&&g^{ab}&=\begin{pmatrix}
1&0&0 \\ 0&\frac{1}{(x^1)^2}&0 \\ 0&0&1
\end{pmatrix}\numberthis\label{eqn:zylinder:Metrik}
\end{align*}
\subsection*{c)}Bestimmung des Gradienten einer Funktion $f(x^1,x^2,x^3)$:
\begin{align*}
\nabla f&=\vec{g}^a\partial_a f=\begin{pmatrix}
\cos x^2\partial_1f-\sin x^2\frac{\partial_2f}{x^1}\\ \cos x^2\partial_1f+\cos x^2\frac{\partial_2f}{x^1} \\ \partial_3f
\end{pmatrix}\numberthis\label{eqn:zylinder:Gradient}
\end{align*}
\subsection*{d)}Berechung der kinetischen Energie in Abhängigkeit von $\dot{x}^a$:
\begin{align*}
T&=\frac{m}{2}g_{ab}\dot{x}^a\dot{x}^b=\frac{m}{2}\kla\left[ {(\dot{x}^1)^2+(x^1)^2(\dot{x}^2)^2+(\dot{x}^3)^2} \right]\numberthis\label{eqn:zylinder:Kinetische Energie}
\end{align*}
\subsection*{e)}Bestimmung von $\Gamma_{22}^1$:
\begin{align*}
\Gamma_{22}^1&=\vec{g}^1\partial_2\vec{g}_2=-x^1\\
\Gamma_{22}^1&=\frac{1}{2}g^{c1}\kla\left( {g_{c2,2}+g_{2c,2}-g_{22,c}} \right)=-\frac{1}{2}g^{11}\partial_1(x^1)^2=-x^1
\end{align*}
\subsection*{f)}


\end{document}